In this chapter, we use some of the rules found for the Broad dataset to predict PD-L1 expression
(specifically whether it exceeds a threshold) in the Sanger dataset.
The performance of these rulesets is also compared to the corresponding optimal ruleset for the Sanger dataset itself.

There are two ways of evaluating a ruleset produced from a train dataset on a test dataset:
\begin{itemize}
    \item \textbf{Absolute inference:} The ruleset is characterized by the absolute values of the thresholds.
    For instance, if the clause is \[\text{STAT2 $\geq -0.20 \; Q(0.5)$},\]
    then \[\text{STAT2 $\geq -0.20$}\] is used for the test dataset.
    \item \textbf{Relative inference:} The ruleset is characterized by the quantile level
    with respect to the observed TF activity values in the inference dataset.
    For instance if the clause is \[\text{STAT2 $\geq -0.20 \; Q(0.5)$},\]
    then \[\text{STAT2 $\geq a$}\] is used for the test dataset,
    where $a$ is the $0.5$ quantile of STAT2 activity in the test dataset.
\end{itemize}
During cross-validation experiments, absolute inference was used, as both train and test datasets were subsets of
the same dataset.
It may be that Sanger is biased with respect to Broad,
for instance due to a difference in measurement techniques or lab environment.
Therefore, relative inference may be the better choice.
Accuracy for both types of inference is reported in \autoref{tab:holdout_C_1}.
Absolute inference structurally yields a lower score, illustrating that clauses are, in this setting, best
interpreted in terms of their relative threshold, meaning their quantile with respect to observed data.
Relative inference yields an accuracy that is more comparable with train accuracy.

The idea is not to evaluate every result obtained on the Sanger dataset, but only those deemed most promising.
The selection made in \autoref{tab:holdout_C_1} is based on the results in \autoref{sec:threshold-results}.

We compare this with the optimal rulesets for the Sanger dataset.
Ideally, the (relative) test accuracies  in \autoref{tab:holdout_C_1} are close to the train accuracies in \autoref{tab:holdout_optimal_C_1}
If not, then the rulesets from Broad do not generalize well, or there is substantial overfitting when training on Sanger.

The optimal rulesets for Sanger for the same parameter choices are given in \autoref{tab:holdout_optimal_C_1}.
These results are qualitatively different, suggesting that the rulesets found for Broad do not generalize to Sanger.
The optimal rulesets for the Sanger dataset also score several percentage points higher, so we do not consider these scores close.

\begin{adjustwidth}{-3cm}{-3cm}
    \begin{table}
        \bwtable{gen/holdout/tab_C_1}
        \caption{Performance of various one-rule rulesets on both Broad and Sanger.
        Performance for absolute inference is indicated  with ``(Abs.)''
            and performance for relative inference is indicated with ``(Rel.)''}
        \label{tab:holdout_C_1}
    \end{table}
\end{adjustwidth}

\begin{adjustwidth}{-3cm}{-3cm}
    \begin{table}
        \bwtable{gen/holdout/tab_optimal_C_1.tex}
        \caption{Optimal rulesets for the Sanger dataset}
        \label{tab:holdout_optimal_C_1}
    \end{table}
\end{adjustwidth}